\documentclass[10pt]{article}

\usepackage[a4paper,margin=22mm]{geometry}
\usepackage{graphicx}
\usepackage{booktabs}
\usepackage{array}
\usepackage{amsmath}
\usepackage{xcolor}
\usepackage{microtype}
\usepackage{newtxtext,newtxmath}
\usepackage{caption}
\usepackage{hyperref}

\definecolor{reviewerblue}{RGB}{0,72,130}
\definecolor{responseblue}{RGB}{0,86,125}
\hypersetup{colorlinks=true,linkcolor=responseblue,urlcolor=responseblue}

\captionsetup{
  font=small,
  labelfont=bf,
  justification=justified,
  singlelinecheck=false
}

\newcommand{\reviewercomment}[1]{%
  \noindent\textbf{\textcolor{reviewerblue}{Reviewer comment:}}%
  \ \textit{#1}\par
}
\newcommand{\authorresponse}{%
  \noindent\textbf{\textcolor{responseblue}{Response:}}\ %
}

\begin{document}

\section*{Response to the Reviewer: Failure-Case Analysis}

\reviewercomment{Notably, one of the recommendations is to include
failure-case analysis for unsuccessful interception and delayed cooperative
engagement scenarios because understanding framework limitations remains
essential for practical battlefield deployment reliability.}

\authorresponse
Thank you for this important recommendation. In the revised manuscript, we
have added a dedicated failure-case analysis for both Case~1 and Case~2. The
purpose of this supplementary experiment is to expose a practically relevant
boundary between \emph{mission-level target coverage} and \emph{strict
completion by every assigned interceptor}. Accordingly, the reported cases are
not catastrophic failures in which attackers penetrate the defense. Instead,
all eight attacking vehicles are neutralized, while one of the twenty assigned
interceptors fails to complete its individual engagement.

\paragraph{Experimental protocol.}
The trained policies, target assignment, attacker trajectories, and the
three-dimensional hit radius \(d_{\mathrm{hit}}=3~\mathrm{m}\) were kept
fixed. No network retraining, backpropagation, optimizer step, or policy-weight
update was performed. To emulate imperfect battlefield information, the
frozen-policy evaluation included a one-sample observation delay
(\(0.05~\mathrm{s}\)), zero-mean position measurement noise with a
\(3~\mathrm{m}\) standard deviation, and zero-mean velocity measurement noise
with a \(0.3~\mathrm{m/s}\) standard deviation. The maximum evaluation horizon
was \(75~\mathrm{s}\), and the cooperative-time tolerance remained
\(0.5~\mathrm{s}\).

We screened 100 reproducible evaluation episodes for each case. To avoid
subjective, appearance-based selection, we used the following fixed rule: the
reported examples are the earliest three seeds in each case for which
(i) all \(8/8\) attacking vehicles were intercepted,
(ii) exactly \(19/20\) assigned interceptors hit,
(iii) seven of the eight target-assignment groups were complete, and
(iv) the arrival-time spread among the observed hits did not exceed the
prescribed \(0.5~\mathrm{s}\) tolerance. No simulated state or trajectory
sample was manually changed.

\begin{figure*}[!htbp]
  \centering
  \includegraphics[width=\textwidth]{%
    reviewer_response_assets/failure_cases_case1_case2_v10.pdf}
  \caption{Representative boundary failures under frozen-policy evaluation
  for Case~1 (a)--(c) and Case~2 (d)--(f). Solid curves denote attacker
  trajectories, dotted curves successful interceptors, and the thick dashed
  curve the missed assigned interceptor. Green circles mark successful hit
  events, while the cross marks the missed interceptor's closest-approach
  point. In every episode, \(19/20\) assigned interceptors hit, but assignment
  redundancy still ensures that all \(8/8\) attacking vehicles are
  intercepted.}
  \label{fig:failure_case_response}
\end{figure*}

\begin{table*}[!htbp]
  \centering
  \caption{Quantitative diagnostics of the six representative boundary
  failures. The miss margin is
  \(d_{\min}^{\mathrm{miss}}-d_{\mathrm{hit}}\).}
  \label{tab:failure_case_response}
  \small
  \setlength{\tabcolsep}{3.2pt}
  \renewcommand{\arraystretch}{1.12}
  \begin{tabular}{ccccccccc}
    \toprule
    Case & Seed & Interceptor hits & Target coverage &
    Missed ID & Target ID &
    \(d_{\min}^{\mathrm{miss}}\) (m) &
    Miss margin (m) & Max.\ spread (s) \\
    \midrule
    1 & 76014 & 19/20 & 8/8 & 1  & 21 & 3.161 & 0.161 & 0.050 \\
    1 & 76048 & 19/20 & 8/8 & 5  & 25 & 3.243 & 0.243 & 0.050 \\
    1 & 76052 & 19/20 & 8/8 & 8  & 20 & 3.041 & 0.041 & 0.050 \\
    2 & 77008 & 19/20 & 8/8 & 1  & 21 & 3.574 & 0.574 & 0.350 \\
    2 & 77020 & 19/20 & 8/8 & 2  & 22 & 3.096 & 0.096 & 0.150 \\
    2 & 77023 & 19/20 & 8/8 & 17 & 21 & 3.488 & 0.488 & 0.200 \\
    \bottomrule
  \end{tabular}
\end{table*}

\paragraph{Observed failure mechanism.}
As shown in Fig.~\ref{fig:failure_case_response} and
Table~\ref{tab:failure_case_response}, assignment redundancy preserved complete
target coverage in all six episodes, although one assigned interceptor did not
enter the lethal region. In Case~1, the missed closest-approach distances were
\(3.041\)--\(3.243~\mathrm{m}\), corresponding to margins of only
\(0.041\)--\(0.243~\mathrm{m}\) above the hit threshold. In Case~2, the
corresponding distances were \(3.096\)--\(3.574~\mathrm{m}\), with margins of
\(0.096\)--\(0.574~\mathrm{m}\). These are therefore near-threshold misses
rather than grossly divergent trajectories.

The seven completed assignment groups remained well coordinated. Their
maximum observed hit-time spread was \(0.05~\mathrm{s}\) in Case~1 and
\(0.15\)--\(0.35~\mathrm{s}\) in Case~2, all below the prescribed
\(0.5~\mathrm{s}\) tolerance. However, after the missed interceptor passed its
closest-approach point, its distance to the assigned target increased and no
hit occurred before the \(75~\mathrm{s}\) horizon. Therefore, the incomplete
group's all-member completion time is right-censored beyond the evaluation
horizon. The same episode consequently represents both an unsuccessful
individual interception and, from the strict group-completion viewpoint, a
delayed cooperative engagement.

\paragraph{Implications for practical reliability.}
This analysis reveals both a strength and a limitation of the proposed
framework. Its redundant assignment structure isolates a single
near-threshold interceptor miss and still neutralizes every attacker, thereby
preventing an individual failure from becoming a mission-level penetration.
Nevertheless, the current framework does not contain an explicit online
recovery mechanism after an interceptor passes its closest-approach point
under delayed and noisy observations. Strict all-member completion therefore
cannot be guaranteed in this boundary regime, even if another member has
already neutralized the assigned target.

We have explicitly added this limitation to the revised manuscript and
identified several practical extensions: (i) online miss detection using
distance rebound and closing-rate reversal, (ii) event-triggered target
reassignment, (iii) activation of a reserve interceptor, and
(iv) uncertainty-aware safety constraints or robust policy training. We also
clarify that the six reported episodes are diagnostic boundary cases intended
to explain the failure mechanism; they are not presented as an estimate of the
overall failure probability. The aggregate reliability of the framework
continues to be evaluated using the Monte Carlo results reported in the main
simulation section.

\end{document}
